\documentclass[11pt]{article}

\usepackage[T1]{fontenc}
\usepackage{amsmath, amssymb}
\usepackage[margin=1in]{geometry}

\title{Introduction to Algorithms (CLRS)\\Chapter 1: The Role of Algorithms}
\author{}
\date{}

\begin{document}
\maketitle

\section{Exercises}

\subsection*{Exercise 1.1-1: Real-world example requiring a convex hull}

A marketing team could plot the locations of people who engaged with a social
media campaign and compute their convex hull. The resulting boundary gives the
team a rough geographic zone around those users to target with a follow-up
campaign. It is an approximation: the hull may include places where there are
no engaged users.

\subsection*{Exercise 1.1-2: Other measures of efficiency}

Besides running time, an algorithm can be judged by how much memory and storage
it uses, how much hardware it requires, and how much energy or network bandwidth
it consumes. These resources have real costs: an algorithm that needs more RAM
may require more expensive machines, while one that takes longer to run uses
more computing time and can increase operating costs. The most important measure
depends on the setting and its constraints.

\subsection*{Exercise 1.1-3: Strengths and limitations of a data structure}

I like trees and graphs because they provide useful abstractions for modeling
relationships and can make solutions easier to adapt to different problems. A
binary search tree, for example, keeps values ordered and can support efficient
searching, insertion, and deletion. Its limitation is that it can become
unbalanced: operations may then take linear time. Keeping it balanced avoids
that degradation, but adds implementation complexity and the cost of
rebalancing after updates.

\subsection*{Exercise 1.1-4: Shortest path and traveling salesman}

Both problems can be represented as weighted graphs and seek a route with a
small total edge weight. In a shortest-path problem, the route connects a
specified start vertex to a specified destination; it does not have to visit
every vertex or return to the start. In the traveling-salesman problem, the
route must visit every vertex (city) and return to its starting point. Finding
shortest paths between individual pairs of cities does not generally solve the
traveling-salesman problem, because the tour must satisfy the visit-every-city
constraint while minimizing its total length.

Their computational difficulty also differs. Shortest-path problems have
polynomial-time algorithms, while the traveling-salesman optimization problem
is NP-hard; its decision version (whether a tour of at most a given length
exists) is NP-complete. Brute force can solve either problem on small inputs,
but checking every possible tour takes factorial time, which quickly becomes
impractical. NP-completeness does not mean brute force is impossible; it means
no polynomial-time algorithm is known for the decision problem.

\subsection*{Exercise 1.1-5: Exact and approximate solutions}

\paragraph{Exact solution.}
A mail carrier needs to deliver letters to every block in a small condominium
and return to the mailroom before a strict cutoff. Model the blocks and walkways
as a weighted graph, then find the shortest tour that visits every block and
returns to the start. With only a few blocks, finding the exact best route can
be practical and shows whether the deadline can be met.

\paragraph{Approximate solution.}
In a very large condominium, there may be too many possible visiting orders to
search for the exact shortest tour quickly. The carrier could instead follow a
simple block-by-block order or repeatedly visit the nearest unvisited block.
The route may involve some extra walking, but it is quick to plan and still gets
every letter delivered. This is a useful trade-off when a near-best route is
sufficient and timely planning matters more than proving the route is optimal.

\subsection*{Exercise 1.2-1: Algorithmic content at the application level}

A web search engine such as Google relies on algorithms to find pages relevant
to a query, rank them, and order the results. These algorithms are central to
the application because their quality directly affects how useful the search
results are.

\subsection*{Exercise 1.2-2: Compare insertion sort and merge sort}

Insertion sort is faster when its step count is smaller:
\begin{equation*}
    8n^2 < 64n \lg n.
\end{equation*}

For positive $n$, divide by $8n$ to get $n < 8\lg n$, where $\lg$ is the
base-2 logarithm. For integer input sizes, this holds for $2 \le n \le 43$.
Starting at $n=44$, merge sort is faster. So insertion sort wins for some sizes
above 8, but not for all sizes above 8.

To verify the cutoff, consider $8\lg n/n$, which must be greater than 1. This
ratio decreases for $n \ge 3$. At $n=43$, it is about $1.0095$, while at
$n=44$, it is about $0.9926$. Thus the inequality stops holding between 43 and
44; check $n=2$ separately, where it also holds.

\subsection*{Exercise 1.2-3: Compare quadratic and exponential running times}

We need the smallest integer $n$ for which
\begin{equation*}
    100n^2 < 2^n.
\end{equation*}

Compare both expressions at the same input size. At $n=14$, they are
$100(14)^2=19{,}600$ and $2^{14}=16{,}384$, so the quadratic algorithm is
still faster. At $n=15$, they are $100(15)^2=22{,}500$ and
$2^{15}=32{,}768$, so the exponential algorithm is faster. The ratio
$2^n/n^2$ increases for integers $n \ge 3$, and at $n=14$ it is still below
100; checking $n=1,2$ separately, the smallest value that satisfies the
inequality is therefore $n=15$.

\end{document}
